\section{Finite-range interactions and general domains}
\label{sec:peps_exact_finite_range}

% Generalizations made in this library. The manuscript states only the
% nearest-neighbor case on the open square; nothing below asserts either
% uniform approximation target.

The finite-size argument of Section~\ref{sec:peps_exact_finite_exception}
uses only the unit norm of the ground vector, and the spanning-tree
construction of Section~\ref{sec:peps_exact_tree} uses only connectivity.
Both therefore extend from the nearest-neighbor square model to the
interaction classes and domains of the area-law
model~\cite[Section~1]{OpenAI2026AreaLaw}. The
manuscript~\cite{OpenAI2026PolynomialPEPS} states only the nearest-neighbor
case; the finite-range target below is a generalization made here, and
neither target is proved.

\begin{theorem}[Exact PEPS on connected lattice domains]
    \label{thm:peps_exact_domain}
    \lean{TNLean.PEPS.AreaLaw.exists_exact_domain_tensor,
        TNLean.PEPS.AreaLaw.domainGraph_rectangularDomain_connected}
    \leanok
    \uses{def:area_law_domain}
    Let $\Lambda\subseteq\mathbb Z^2$ be a finite domain whose induced
    nearest-neighbor graph is connected, and let $q\ge1$. Every
    $\psi\colon\{0,\ldots,q-1\}^{\Lambda}\to\mathbb C$ is the exact
    contraction of a PEPS on the induced graph with bond dimensions
    $1\le D_e\le q^{|\Lambda|}$ on every edge. The domain of a $w\times h$
    rectangle with $w,h>0$ is connected.
\end{theorem}
\begin{proof}
    \leanok
    \uses{thm:peps_exact_degenerate,thm:peps_square_connected,thm:area_law_domain_identifications}
    Theorem~\ref{thm:peps_exact_degenerate} applies to the induced graph,
    whose vertex set has $|\Lambda|$ elements. The rectangular domain is
    isomorphic to the open rectangular grid
    (Theorem~\ref{thm:area_law_domain_identifications}), which is connected
    by Theorem~\ref{thm:peps_square_connected}, and connectivity is
    preserved by graph isomorphisms.
\end{proof}

For $R\in\mathbb N$, real $J,\Delta$ and $L\in\mathbb N$, let
$\Lambda_L=\{0,\ldots,L-1\}^2$ be the integer domain of the open $L\times L$
square and let $U$ identify the configuration basis of the square with that
of $\Lambda_L$ (Theorem~\ref{thm:area_law_domain_identifications}). Write
$\mathcal G^{R}_{q,J,\Delta}(L)$ for the triples $(H,E_0,\Omega)$ in which
$H$ is an interaction family on $\Lambda_L$ as in
Definition~\ref{def:area_law_hamiltonian}, of range $R$ and term bound $J$,
and $\Omega$ is a vector on the open square such that $U\Omega$ is a unit
vector with
\begin{align}
    HU\Omega & =E_0U\Omega,
             & H-E_0I       & \ge\Delta\bigl(I-|U\Omega\rangle\langle U\Omega|\bigr).
    \label{eq:peps_finite_range_ground}
\end{align}
By Theorem~\ref{thm:rectangular_domain_hamiltonian}, for $J\ge0$ every
$(H,E_0,\Omega)\in\mathcal G_{q,J,\Delta}(L)$ yields a member of
$\mathcal G^{1}_{q,J,\Delta}(L)$ with the same $E_0$ and $\Omega$.

\begin{theorem}[A uniform prefactor for finite-range interactions]
    \label{thm:peps_remove_finite_exceptions_finite_range}
    \lean{TNLean.PEPS.AreaLaw.exists_uniform_finiteRange_approximation_of_sufficiently_large}
    \leanok
    \uses{def:area_law_hamiltonian,thm:rectangular_domain_hamiltonian}
    Fix integers $q\ge1$, $R\ge0$ and $L_0\ge0$, real $J,\Delta,C$, and
    $c\ge0$. Suppose $\mathcal A_{C,c}(L,\Omega)$ holds for every integer
    $L\ge\max\{2,L_0\}$ and every
    $(H,E_0,\Omega)\in\mathcal G^{R}_{q,J,\Delta}(L)$. There is one $C'>0$
    with $C'\ge C$ such that $\mathcal A_{C',c}(L,\Omega)$ holds for every
    integer $L\ge2$ and every $(H,E_0,\Omega)\in\mathcal G^{R}_{q,J,\Delta}(L)$.
    The exponent $c$ is unchanged, and $C'$ is independent of
    $L,H,E_0,\Omega$.
\end{theorem}
\begin{proof}
    \leanok
    \uses{thm:peps_remove_finite_exceptions_unit}
    Apply Theorem~\ref{thm:peps_remove_finite_exceptions_unit} to the sets
    $\mathcal P(L)$ of vectors $\Omega$ occurring in
    $\mathcal G^{R}_{q,J,\Delta}(L)$. Their members are unit vectors because
    $U$ preserves the norm.
\end{proof}

\begin{definition}[Finite-range uniform polynomial approximation]
    \label{def:peps_finite_range_uniform_target}
    \lean{TNLean.PEPS.AreaLaw.FiniteRangePolynomialPEPSApproximation}
    \leanok
    \uses{def:area_law_hamiltonian,thm:rectangular_domain_hamiltonian,def:polynomial_peps_model}
    The finite-range uniform polynomial approximation assertion is
    \begin{align}
         & \forall q\ge2\;\forall R\in\mathbb N\;\forall J,\Delta>0\;
        \exists C,c>0\;\label{eq:peps_finite_range_uniform}           \\
         & \quad\forall L\ge2\;
        \forall(H,E_0,\Omega)\in\mathcal G^{R}_{q,J,\Delta}(L),\;
        \mathcal A_{C,c}(L,\Omega).\notag
    \end{align}
    The constants $C,c$ depend only on $q,R,J,\Delta$. This generalizes
    assertion~(\ref{eq:peps_finite_uniform}) beyond the nearest-neighbor
    interactions of~\cite[Theorem~1.1]{OpenAI2026PolynomialPEPS}; it is not
    asserted here.
\end{definition}

\begin{theorem}[Equivalence with the sufficiently large size statement]
    \label{thm:peps_finite_range_approximation_iff_large}
    \lean{TNLean.PEPS.AreaLaw.finiteRangePolynomialPEPSApproximation_iff_sufficiently_large}
    \leanok
    \uses{def:peps_finite_range_uniform_target}
    Assertion~(\ref{eq:peps_finite_range_uniform}) is equivalent to
    \begin{align}
         & \forall q\ge2\;\forall R\in\mathbb N\;\forall J,\Delta>0\;
        \exists C,c>0\;\exists L_0\in\mathbb N\;
        \label{eq:peps_finite_range_large}                            \\
         & \quad\forall L\ge\max\{2,L_0\}\;
        \forall(H,E_0,\Omega)\in\mathcal G^{R}_{q,J,\Delta}(L),\;
        \mathcal A_{C,c}(L,\Omega).\notag
    \end{align}
    In both assertions the constants are independent of the size,
    interaction family, ground energy, and ground vector.
\end{theorem}
\begin{proof}
    \leanok
    \uses{thm:peps_remove_finite_exceptions_finite_range}
    As for Theorem~\ref{thm:peps_approximation_iff_large}:
    assertion~(\ref{eq:peps_finite_range_uniform}) gives
    (\ref{eq:peps_finite_range_large}) with $L_0=0$, and conversely
    Theorem~\ref{thm:peps_remove_finite_exceptions_finite_range} produces one
    positive $C'\ge C$ for every $L\ge2$ with the same exponent.
\end{proof}

\begin{theorem}[Specialization to nearest-neighbor interactions]
    \label{thm:peps_finite_range_implies_nearest_neighbor}
    \lean{TNLean.PEPS.AreaLaw.FiniteRangePolynomialPEPSApproximation.polynomialPEPSApproximation}
    \leanok
    \uses{def:peps_finite_range_uniform_target,def:polynomial_peps_model}
    Assertion~(\ref{eq:peps_finite_range_uniform}) implies the
    nearest-neighbor assertion~(\ref{eq:peps_finite_uniform}).
\end{theorem}
\begin{proof}
    \leanok
    \uses{thm:rectangular_domain_hamiltonian,thm:rectangular_equal_side_entropy}
    Fix $q\ge2$ and $J,\Delta>0$, and take the constants $C,c$
    of~(\ref{eq:peps_finite_range_uniform}) at $R=1$. For
    $(H,E_0,\Omega)\in\mathcal G_{q,J,\Delta}(L)$, the square Hamiltonian is
    the equal-side case of the rectangular model, and
    Theorem~\ref{thm:rectangular_domain_hamiltonian} converts it into a
    range-one family $H'$ on $\Lambda_L$ with ground vector $U\Omega$, energy
    $E_0$ and gap $\Delta$. Thus
    $(H',E_0,\Omega)\in\mathcal G^{1}_{q,J,\Delta}(L)$, and
    $\mathcal A_{C,c}(L,\Omega)$ follows.
\end{proof}
